%! suppress = LineBreak
%% For double-blind review submission, w/o CCS and ACM Reference (max submission space)
%\documentclass[sigplan,10pt,review,anonymous]{acmart}
%\settopmatter{printfolios=false,printccs=false,printacmref=false}
%% For double-blind review submission, w/ CCS and ACM Reference
%\documentclass[sigplan,review,anonymous]{acmart}\settopmatter{printfolios=true}
%% For single-blind review submission, w/o CCS and ACM Reference (max submission space)
%\documentclass[sigplan,review]{acmart}\settopmatter{printfolios=true,printccs=false,printacmref=false}
%% For single-blind review submission, w/ CCS and ACM Reference
%\documentclass[sigplan,review]{acmart}\settopmatter{printfolios=true}
%% For final camera-ready submission, w/ required CCS and ACM Reference
%\documentclass[sigplan,nonacm]{acmart}
\documentclass[sigplan,acmsmall,nonacm,review,anonymous]{acmart}%\settopmatter{printfolios=true,printccs=true,printacmref=true}

%% Conference information
%% Supplied to authors by publisher for camera-ready submission;
%% use defaults for review submission.
\acmConference[POPL '27]{The 54th Annual ACM SIGPLAN Symposium on Principles of Programming Languages}{January 10--16, 2027}{Mexico City, Mexico}
\acmConference{}{}{}
\acmYear{2027}
\acmISBN{} % \acmISBN{978-x-xxxx-xxxx-x/YY/MM}
\acmDOI{} % \acmDOI{10.1145/nnnnnnn.nnnnnnn}
\startPage{1}

%% Copyright information
%% Supplied to authors (based on authors' rights management selection;
%% see authors.acm.org) by publisher for camera-ready submission;
%% use 'none' for review submission.
%\setcopyright{none}
%\setcopyright{acmcopyright}
\setcopyright{acmlicensed}
%\setcopyright{rightsretained}
%\copyrightyear{2018}           %% If different from \acmYear

%% Bibliography style
\bibliographystyle{acmart}

\input{preamble}

%\usepackage{draftwatermark}
%\SetWatermarkLightness{0.75}
%\SetWatermarkText{DRAFT}
\makeatletter
\let\@authorsaddresses\@empty
\makeatother

% arXiv does not support LuaLaTeX. Use the following options to submit:
%\pdfextension info {/Creator() /Producer()}
%\pdfvariable suppressoptionalinfo \numexpr 1+2+4+8+16+32+64+128+256+512 \relax
% Recompile, then run: sed -e '/PTEX\./s/./ /g' < popl.pdf > tidyparse.pdf
% Finally, use "Print to PDF" on MacOS Preview.

\begin{document}
%
\title{Syntax Repair as Language Intersection}
%
\begin{abstract}
We introduce a new technique for repairing syntax errors in arbitrary context-free languages. This technique models syntax repair as a language intersection problem by defining a finite language that provably generates every syntactically valid repair within a given edit distance. Leveraging a theoretical connection between the Bar-Hillel construction from formal language theory and CFL reachability from program analysis, we give a parallel matrix-closure procedure for deciding repairability in a finite number of typographic edits and provide an enumeration algorithm based on the Brzozowski derivative. Finally, we evaluate this algorithm and its implementation, demonstrating state-of-the-art results on a Python syntax repair benchmark.\keywords{Error correction \and CFL reachability \and Language games.}
\end{abstract}

%\titlerunning{Abbreviated paper title}
% If the paper title is too long for the running head, you can set
% an abbreviated paper title here
\author{Breandan Considine}
\email{bre@ndan.co}

\maketitle

\subfile{sec-introduction}
\subfile{sec-background}
\subfile{sec-method}
\subfile{sec-examples}
\subfile{sec-measuring-intersection}
\subfile{sec-decoding-reranking}
\subfile{sec-evaluation}
\subfile{sec-discussion}
\subfile{sec-related-work}
\subfile{sec-conclusion}%
\clearpage\bibliography{../bib/acmart}\vspace{-1cm}

\pagebreak\appendix

\subfile{sec-implementation}
\subfile{app-levenshtein-topology}
\subfile{app-levenshtein-minimality}
\subfile{app-hyperparameters}
\subfile{app-example-repairs}
\subfile{app-raw-data}
\subfile{app-symbols}
\clearpage\subfile{app-chatgpt-repair-prompt}%
\end{document}
